\section{第4组：流体力学}
\sourceproblem{4-T1}{13}{17}
\problemtext 水平变截面管稳定流动，出口面积为最细处三倍，出口速率2米每秒，求最细处速度及表压，并判断细处小孔出水情形。
\workings 连续性给出 $v_1=6\,\mathrm{m/s}$。出口通大气，忽略损失，伯努利方程
\[
p_1-p_{\rm atm}=\tfrac12\rho(2^2-6^2)=-16\rho\,\mathrm{m^2/s^2}.
\]
取水密度1000千克每立方米，表压为 $-1.6\times10^4\,\mathrm{Pa}$，绝对压强为 $p_{\rm atm}-1.6\times10^4\,\mathrm{Pa}$；小孔处通常吸入空气而非向外喷水。
\crosscheck 结论以孔未开启前稳态、出口通大气且无损为前提；开启后原有流场可能改变。
\sourceproblem{4-T2}{13}{17}
\problemtext 水平管A面积10平方厘米，分支B、C面积8、6平方厘米，$v_A=1.00,v_C=0.50$（m/s），求B速度、B与C压差及最大压强所在管。
\workings $A_Av_A=A_Bv_B+A_Cv_C$，故 $v_B=0.875\,\mathrm{m/s}$。理想等高支路共享总压，
\[
p_C-p_B=\tfrac12\rho(v_B^2-v_C^2)\approx258\,\mathrm{Pa}.
\]
C速度最小，故C压强最大。
\crosscheck 假设无分支局部损失且各支路有同一上游总压；实际管网不能仅凭流速确定所有压差。
\sourceproblem{4-T3}{14}{18}
\problemtext 开口水槽液面高 $H$，孔在液面下深度 $h$，求水平射程、另一等射程孔位置、最大射程。
\workings 理想小孔流速 $\sqrt{2gh}$；孔距地 $H-h$，下落时间 $\sqrt{2(H-h)/g}$。水平射程
\[
X=2\sqrt{h(H-h)}.
\]
交换 $h$ 与 $H-h$ 射程不变，故另一孔深度 $H-h$；当 $h=H/2$ 时最大，$X_{\max}=H$。
\crosscheck $0<h<H$；两孔关于半高平面对称；水面下降及喷流阻力忽略。
\sourceproblem{4-T4}{14}{18}
\problemtext 顶部开口直立圆柱水箱，初水深0.1米，底孔面积为横截面积的1/400，求理想排空时间。
\workings 设水深 $h$，面积比 $A/a=400$，连续性及托里拆利公式给出 $A\dot h=-a\sqrt{2gh}$，积分
\[
t=\frac{A}{a}\sqrt{\frac{2h_0}{g}}
=400\sqrt{\frac{0.2}{g}}\approx57.1\,\mathrm s
\]
（$g=9.8$）。若保留自由液面速度，乘以 $\sqrt{1-(a/A)^2}$，变化极小。
\crosscheck 液面高度变化不可用初始出流速度直接除以体积；真实排水需考虑流量系数。
